\section{Troubleshooting}
\label{sec:Troubleshooting}

This chapter covers common failure modes, diagnostic approaches and recovery
procedures.

\subsection{Diagnostic endpoints}

Several endpoints help diagnose problems without restarting the server:

\begin{longtable}{p{54mm}p{100mm}}
\toprule
Endpoint & Purpose \\
\midrule
\texttt{GET /admin/overview} & Check which models exist and their health.
  Look for stale head revisions or integrity warnings. \\
\texttt{GET /admin/models/\{modelId\}/status} & Detailed model status:
  counts, revision, consistency. \\
\texttt{GET /admin/models/\{modelId\}/integrity} & Structural integrity:
  missing references, orphan nodes. \\
\texttt{GET /admin/auth/diagnostics} & Authorization troubleshooting:
  resolved identity, roles, mode. \\
\texttt{GET /admin/sessions} & Active WebSocket sessions across all
  models. Check for stale or abandoned sessions. \\
\texttt{GET /admin/audit/config} & Current audit configuration. \\
\bottomrule
\end{longtable}

\subsection{Common problems}

\subsubsection{Client cannot connect}

\begin{description}
\item[Symptom] Connection fails with "Failed to connect archimesh websocket".
\item[Possible causes]
  \begin{itemize}
  \item Server is not running. Check with \\
    \texttt{curl http://localhost:8081/admin/overview}.
  \item Wrong URL. The WebSocket URL should be \\
    \texttt{ws://<host>:<port>} (not \texttt{http}).
  \item Authorization failure. Check the auth token validity and role
    assignments. Use \texttt{GET /admin/auth/diagnostics}.
  \item Model does not exist. Models must be registered via
    \texttt{POST /admin/models/\{modelId\}} before clients can connect.
  \item Network firewall blocking WebSocket upgrade.
  \end{itemize}
\end{description}

\subsubsection{SubmitOps rejected}

\begin{description}
\item[Symptom] The server broadcasts \texttt{Error} messages to the client.
  The status line shows a conflict banner.
\item[Error codes]
  \begin{description}
  \item[\texttt{REVISION\_AHEAD}] Client's \texttt{baseRevision} exceeds
    server head. The client is ahead of the server --- this usually means
    the client missed a reconnect after a server state change. Solution: the
    client should disconnect and reconnect (the server signals this by
    broadcasting the error to all sessions).
  \item[\texttt{PRECONDITION\_FAILED}] A referenced entity does not exist.
    The op references an element, relationship, view, view object or
    connection that is not in the materialized graph and was not created
    earlier in the same batch. Common causes: deleting an entity on one
    client while another client edits a diagram object representing it;
    submitting a connection create before the relationship is committed.
    The conflicting outbox entry is dropped and surfaced in the UI.
  \item[\texttt{LOCK\_CONFLICT}] Another client holds a notation lock on
    one of the targets. The client should retry after releasing and
    re-acquiring the lock. The conflicting outbox entry is dropped.
  \end{description}
\end{description}

\subsubsection{Materialized state corrupted}

\begin{description}
\item[Symptom] Integrity check reports missing references or orphan nodes.
\item[Cause] Bug in the apply path, manual Neo4j modification or failed
  partial writes.
\item[Solution] Run a rebuild: \\
  \texttt{curl -s -X POST http://localhost:8081/models/demo/rebuild}
\item[Verification] After rebuild, run integrity check again and verify
  \texttt{issues} is empty.
\end{description}

\subsubsection{Neo4j connection failure}

\begin{description}
\item[Symptom] Server logs show Neo4j connection errors. All operations fail.
\item[Check] Is Neo4j running? \texttt{docker compose ps neo4j}. Check the
  Neo4j browser at \url{http://localhost:7474}.
\item[Check] Are credentials correct? Verify \texttt{NEO4J\_URI},
  \texttt{NEO4J\_USER} and \texttt{NEO4J\_PASSWORD} match the running
  instance.
\item[Check] Is the Bolt port accessible? Default is 7687.
\end{description}

\subsubsection{Kafka unavailable}

\begin{description}
\item[Symptom] Server logs show Kafka publish warnings. Other clients do not
  receive real-time updates from the submitter's client.
\item[Impact] Submitting clients still receive \texttt{OpsAccepted} and can
  continue editing. The server falls back to direct WebSocket broadcast for
  \texttt{OpsBroadcast}, so clients on the same server instance still
  receive updates.
\item[Check] Is Kafka running? \texttt{docker compose ps kafka}. Check the
  Kafka UI at \url{http://localhost:8080}.
\item[Check] Do the required topics exist? Use \\
  \texttt{docker compose exec kafka rpk topic list --brokers localhost:9092}.
\end{description}

\subsubsection{Outbox filling up}

\begin{description}
\item[Symptom] The Archi status line shows "Outbox near capacity". Local
  edits are being queued because the client is disconnected or the
  connection is flapping.
\item[Impact] When the outbox reaches 1000 entries, the oldest entries are
  silently dropped. The warning fires at 800 entries (80\%).
\item[Solution] Reconnect to the server so the outbox can drain. If
  connectivity is persistently poor, work in smaller edit batches and
  reconnect frequently.
\item[Prevention] Monitor the server's session page for clients that
  repeatedly connect and disconnect. Investigate network stability.
\end{description}

\subsubsection{Cannot overwrite a model}

\begin{description}
\item[Symptom] Import with \texttt{?overwrite=true} returns 409.
\item[Cause] Active sessions exist on the target model.
\item[Solution] Disconnect all clients from the model, verify with
  \texttt{GET /admin/sessions}, then retry. If a session is stuck, force
  delete the model and re-import.
\end{description}

\subsection{Log analysis}

The server emits structured logs. Key log categories:

\begin{longtable}{p{46mm}p{108mm}}
\toprule
Logger & Content \\
\midrule
\texttt{org.gautelis.archimesh} & Application-level events: submit, accept,
  reject, join, disconnect, rebuild, compact, export, import, tags. \\
\texttt{org.apache.kafka} & Kafka client logs (WARN by default). \\
\texttt{admin\_audit} (INFO) & Admin API access audit events. \\
\texttt{ws\_audit} (INFO) & WebSocket lifecycle audit events. \\
\bottomrule
\end{longtable}

Set \texttt{ARCHIMESH\_LOG\_LEVEL=DEBUG} for detailed flow tracing including
op-level detail. Note that DEBUG logging of \texttt{OpsBroadcast} messages
can produce very large log volumes on active models.

Log format:

\begin{lstlisting}[caption={Log format}]
yyyy-MM-dd HH:mm:ss.SSS LEVEL [logger.] (thread) message
\end{lstlisting}

Audit events include a stable JSON payload with user id, action and request
context. Parse the JSON rather than depending on free-form message text.

\subsubsection{Client-side log analysis}

The client plugin writes to the Archi workspace log
(\filepath{\textasciitilde/Archi/workspace/.metadata/.log}). Search for:

\begin{itemize}
\item \texttt{archimesh} --- all plugin log entries
\item \texttt{Connected archimesh} --- successful connections
\item \texttt{SubmitOps received} --- server acknowledging submits
\item \texttt{Queued SubmitOps} --- entries moved to outbox
\item \texttt{Outbox full} --- outbox dropping entries
\item \texttt{Dropping stale} --- conflict resolution dropping an entry
\item \texttt{Error stopping Archimesh} --- shutdown errors
\end{itemize}

The \texttt{scripts/check-archimesh-sanity.sh} script scans Archi workspace
logs for critical sync issues, integrity warnings and connection errors.

\subsection{Recovery procedures}

\subsubsection{Server restart}

If the server is restarted:

\begin{enumerate}
\item All WebSocket connections are closed. Clients will show
  "Disconnected" in the status line.
\item Client outboxes start queuing local edits.
\item When the server is back up, clients must manually reconnect.
\item On reconnect, the client sends a \texttt{Join} with its cached
  revision. The server sends a delta (if the revision gap is within the
  op-log window) or a full snapshot.
\item The outbox drains automatically after reconnect and revision hints
  are available.
\end{enumerate}

\subsubsection{Neo4j data loss}

If the Neo4j database is lost or irrecoverably corrupted:

\begin{enumerate}
\item Restore from the most recent export package using
  \texttt{POST /admin/models/import}.
\item If no export exists, the model timeline is lost. Re-register the
  model and let clients push their local state as new submit batches.
\item Clients with cached models can reconnect after server restore;
  however, the server will have no op-log history, so the client must
  discard its cache and cold-start with a snapshot from the server (which
  will be empty for a newly registered model).
\end{enumerate}

\subsubsection{Full system recovery}

In a disaster scenario where both Neo4j and the server are lost:

\begin{enumerate}
\item Restore infrastructure: \texttt{docker compose up -d}.
\item Apply Neo4j schema: \texttt{docker compose exec -T neo4j cypher-shell ... < neo4j/schema.cypher}.
\item Start the server.
\item Import the most recent model exports.
\item Clients reconnect and receive snapshots from the restored server.
\end{enumerate}

\subsection{Performance tuning}

\begin{itemize}
\item \textbf{Neo4j heap:} The materialized graph and op-log grow with model
  activity. Monitor Neo4j heap usage and increase \texttt{NEO4J\_dbms\_memory\_heap\_max\_\_size}
  as needed.
\item \textbf{Kafka retention:} Long Kafka retention on ops topics consumes
  disk. If rebuild is done from Neo4j (not Kafka), set shorter retention
  for ops topics.
\item \textbf{Compaction frequency:} Run compaction periodically on active
  models to bound op-log growth. The compaction reports the number of
  pruned commits and remaining eligible metadata.
\item \textbf{Poll interval:} The admin UI polls the server for status
  updates. Set the poll interval to 0 to disable auto-refresh if dashboard
  traffic is a concern.
\item \textbf{Consistency checks:} Keep \texttt{ARCHIMESH\_CONSISTENCY\_CHECK}
  disabled in production unless debugging. It adds latency to every submit.
\end{itemize}
